\subsection{Overview}

Subsurface (tile) drainage is one of the best-established engineering interventions for leaching accumulated salts out of irrigated and waterlogged soil profiles, and the assembled evidence---spanning field experiments, lysimeter and soil-column studies, long-term monitoring, and calibrated numerical models across arid, semi-arid, monsoonal, coastal, and cold-climate settings---converges on a coherent picture of what the technology accomplishes and how. Three headline conclusions are strongly supported. First, installing subsurface drainage reliably lowers the water table below the root zone and removes salt from the upper profile, converting saline and waterlogged land into productive cropland, with topsoil salinity reductions commonly of 35--80\% and corresponding crop-yield gains; a global meta-analysis of 891 experiments confirms that drainage significantly reduces soil salinity and increases yield relative to no drainage \parencite{Yang2025}, and a full production-scale field study validates outcomes previously seen only at small scale \parencite{Derby2024}. Second, the rate and completeness of desalinization are governed by a well-characterized set of design and operational parameters---drain spacing, drain depth, envelope/filter, leaching-water quantity, and complementary practices---such that closer spacing, greater depth (up to soil-specific limits), and larger applied leaching volume shorten the desalinization period. Third, the salt exported in drainage water is initially high and declines over years as the profile is flushed, so that salt-balance accounting (salt removed versus salt added) and the temporal trajectory of effluent salinity are central to both agronomic and environmental evaluation. The remainder of this section synthesizes the consolidated findings on reclamation efficacy, the design parameters that control salt removal, drainage-water salinity and salt balance, the vertical and ionic structure of leaching, the hydraulics of efficient leaching and enhancement techniques, saline/brackish water reuse, integration with complementary practices, and the modeling and monitoring tools that underpin these conclusions.

\subsection{Reclamation Efficacy: Soil Desalinization and Crop Response}

Across widely differing environments, subsurface drainage produces substantial and repeatable reductions in root-zone salinity, typically concentrated in the surface layers, together with improved crop performance. In the Northern Great Plains, topsoil salinity in a tile-drained area declined over time while an adjacent undrained area showed little change or a slight increase; drainage alleviated both salinity and wetness constraints, and soybean and corn yields improved within three years, with continued soybean gains expected over at least eight years \parencite{Derby2024}. Comparable field-scale reclamations are documented in India, where average soil salinity in drained blocks fell from 7.18 to 4.62~dS~m$^{-1}$ (an average 35.7\% decrease, ranging 16.0--66.3\% across blocks) while the undrained area rose from 9.00 to 9.54~dS~m$^{-1}$, and crop yields rose by 18.75--27.62\% for wheat, rice, pearl millet, and sorghum \parencite{Sharma2006}. In monsoonal India, drainage combined with monsoon leaching brought initially crop-limiting salinity down to levels suitable for irrigated agriculture, with weighted-average ECe after 13 monsoon leachings of 1.4, 2.1, and 3.4~dS~m$^{-1}$ at 25, 50, and 75~m spacing \parencite{Sharma2000}. Rolling-topography reclamation under unusually wide drain spacing lowered topsoil salinity from 8.4 to about 2.1~dS~m$^{-1}$ and raised paddy yields, with cropping intensity increasing from 143\% to 191\% \parencite{Manjunatha2004}, and multi-scenario evaluation of the Tungabhadra command reduced head- and middle-reach topsoil ECe from roughly 9~dS~m$^{-1}$ to 1.9--3.8~dS~m$^{-1}$, with weaker effect in the tail reach \parencite{Karegoudar2022}.

Reclamation magnitudes are similarly large in China's arid and coastal saline lands. Subsurface pipe drainage reduced 0--80~cm salinity from 35 to 15~g~kg$^{-1}$, whereas open-ditch controls remained at 25--35~g~kg$^{-1}$ \parencite{Wang2020}; a drip-irrigation-plus-subsurface-pipe system cut soil EC from 14.42--15.74 to 2.86--6.91~dS~m$^{-1}$ over three leaching cycles and raised seed cotton yield 52\% \parencite{Heng2021}; and after four years the salinity of the 0--1.0~m layer fell from 27.6 to 4.1--7.2~dS~m$^{-1}$, reclaiming extremely saline soil to moderately saline \parencite{Zhang2021}. One-year monitoring by fixed electromagnetic-induction profiles showed soil apparent EC falling 40.0--62.0\% across depths (mean 49.5\%) and a sharp reduction in spatial heterogeneity \parencite{Yu2025}. Surface-concentrated desalinization is a recurrent signature: 0--20~cm desalination rates exceeded 80\% while deeper layers lagged or even accumulated salt \parencite{Han2023}, and the meta-analysis found the most substantial reduction at depths of 20~cm or less, diminishing with depth and becoming nonsignificant below 80~cm \parencite{Yang2025}. Reductions of comparable magnitude are reported for Turkey (topsoil salinity from 10.3 to 2.0~dS~m$^{-1}$, about 80\% at 20~cm) \parencite{Bahçeci2009}, Tunisia (mean ECe decrease 23.3\%) \parencite{HajAmor2019}, and semi-arid Canada, where combined drainage and windbreak treatment lowered 0--450~mm ECe from 14.08 to 9.58~dS~m$^{-1}$ and raised barley yield up to twentyfold \parencite{Steppuhn2006,Steppuhn2017}. Spatial reclamation is accompanied by homogenization: leaching plus drainage shifted salinity from strong to moderate variability (coefficient of variation from 48.13\% to 28.14\%) and expanded mildly saline area to 76\%, with 7~g~kg$^{-1}$ identified as the critical salinity for the homogenizing tendency \parencite{Yu2015}.

Improved soil physical condition and productivity accompany desalinization. Drainage increased porosity, available water, and infiltration rate (from 5.3 to 13.6--16.0~mm~h$^{-1}$) while lowering bulk density, and reduced SARe and ESP \parencite{Sharma2000}; the meta-analysis reports increased soil permeability and organic matter and decreased CO$_3^{2-}$, HCO$_3^-$, Cl$^-$, SO$_4^{2-}$, Na$^+$, and Mg$^{2+}$ \parencite{Yang2025}. In Egypt, drainage raised yields by 13--16\% across wheat, rice, cotton, and maize and lifted net farm income by 9--44\% depending on region \parencite{Kotb2000}. Where drainage is absent or the water table stays shallow, combined waterlogging and salinity are more damaging than waterlogging alone, and yields fail above 12~dS~m$^{-1}$ \parencite{Kahlown2002}.

\subsection{Design Parameters Governing Salt Removal}

A consistent, quantitatively supported relationship links subsurface-drainage intensity to desalinization: closer drain spacing and greater drain depth shorten the desalinization period and increase salt removal, though depth benefits saturate at soil-specific limits and shallow drains lose effectiveness. A DRAINMOD-based analysis showed the desalinization period was negatively related to drainage intensity; at 1.2~m depth, salinity fell 55.5\% in the first two years at 20~m spacing but only 5.9\% in five years at 200~m spacing, and reaching a slightly saline target within a decade required spacing below 30~m \parencite{Chen2025}. The meta-analysis generalizes this: salinity reduction improves with smaller spacing, greater depth, and longer operation time, with climate- and groundwater-specific layout recommendations---spacing $>$30~m and depth $>$1.2~m in humid climates, 10--20~m spacing with depth $>$1.2~m in arid climates with water table $\leq$1~m \parencite{Yang2025}. Field and column studies repeatedly identify optimum spacings: 15~m gave the best desalination (effective rate 65.8\%, salt-removal ratio 69.6\%) among 10--25~m spacings \parencite{Wang2020}; 20~m spacing at 1.5~m depth optimized root-zone desalination (39.02\% reduction versus control) by multi-objective optimization \parencite{Wang2025c}; 6~m spacing with 40~cm outlet depth maximized average desalination \parencite{Han2023}; and 30~cm partition spacing outperformed wider spacings \parencite{Chen2022}. Smaller spacing also lowers the leaching quota needed for a given efficiency and deepens and stabilizes the water table \parencite{Yang2022c}, while desalination benefits decay once spacing exceeds about 30~m, with some wide-spacing/deep-burial treatments even showing negative desalination \parencite{Zhu2025b}.

Drain depth effects are pronounced but bounded. Pipes shallower than 1.0~m were essentially invalid for discharging salt or controlling the water table \parencite{Qian2021b}, and desalination and salt discharge indices correlated positively with depth \parencite{Qian2021b,Tian2023c,Tian2024}. Yet SaltMod and SahysMod simulations show diminishing returns beyond a threshold: root-zone salinity decreased with depth up to 1.4~m with no further change beyond it \parencite{Bahçeci2007}, and increasing depth from 0.80 to 1.60~m only slightly reduced salinity while increasing discharge \parencite{Bahçeci2006}. For crop-yield optimization under saline irrigation, an intermediate depth is frequently preferred---0.8~m gave the best leaching and highest maize yield \parencite{Feng2017}, relative yield and water-use efficiency peaked near 90--100~cm \parencite{Feng2019}, and reducing depth from 140 to 100--120~cm saved 10--15\% of irrigation water with only a small yield penalty and no significant long-term salinity increase \parencite{Muhammad2021}. Envelope and filter design is a first-order control on discharge and salt leaching: the envelope factor dominated over pipe material, with spun-bonded non-woven geotextile, single layer, and 3.0\% perforation optimal \parencite{Li2022}; double-layer geotextile integrated filtration, anti-clogging, and drainage best (all scenarios $>$95\% salt removal), while triple layers caused filter-cake clogging \parencite{Wang2024d}; and the O$_{90}$/d$_{90}$ ratio predicts clogging and soil loss into pipes \parencite{Elzoghby2021}. Complex stratigraphy modifies the optimum: low-permeability silty interlayers reduce leaching, and pipes should be placed below such interlayers, with desalination increasing linearly with depth and interlayer conductivity and exponentially as spacing decreases \parencite{Guo2024b}; a shallow red-clay layer likewise reduces drainage and salt-removal efficiency \parencite{Tian2024}. Long-term modeling indicates deeper, wider drains desalinize the whole profile better whereas shallow, close drains act mainly in the top 80~cm \parencite{KAMRA1991}, and economic analyses in multiple settings converge on roughly 20~m spacing at 1.2~m depth as the maximum-return layout \parencite{Chen2025}.

\subsection{Drainage-Water Salinity, Salt Balance, and Temporal Decline}

Because tile drains intercept the salt-laden water flushed from the profile, drainage-water salinity is characteristically much higher than irrigation-water salinity and declines over years as the reservoir of stored salt is depleted. Effluent EC exceeding irrigation EC has been documented repeatedly \parencite{Li2022,Karegoudar2022}, with ECdw/ECiw ratios of 11--17 and annual salt export near 29.4~t~ha$^{-1}$ in a high-percolation Iranian sugarcane district \parencite{Akram2013}, and 8{,}000~mg~L$^{-1}$ irrigation water producing 14{,}000~mg~L$^{-1}$ drainage water through evapotranspirative concentration in California \parencite{KARAJEH1994c}. Salt-balance accounting is the standard diagnostic of reclamation: salt removed exceeded salt added in the first year (125.4 versus 119.8~mg), although over the full record (August 1998--April 2002) cumulative salts added ($\approx$452.9~mg) exceeded salts removed ($\approx$372.9~mg) \parencite{Manjunatha2004}, an oasis system exported about 94\% of the imported salt load while discharging 42\% of irrigation water, yielding a positive salt balance \parencite{Askri2022}, and a straw-drainage module removed about 1.7 times the input salt mass \parencite{Lu2019}. Salt removal can, however, fall short: an oasis-scale system removed only 30\% then 15\% of salts owing to mineral and root clogging \parencite{BenAissa2013}, and a sand-clogged gravel filter drained no more than 27\% of the imported salt \parencite{HajAmor2019}. During the irrigation--drainage season, subsurface pipes discharged 3.12--4.29~t~ha$^{-1}$ of salt, accounting for 34--46\% of soil salt-storage change, the remainder leaching to deeper layers or groundwater \parencite{Liu2025c}; over four years pipes removed 40\% of total salt discharge \parencite{Heng2021}.

The temporal decline of effluent salinity is a robust, long-validated pattern. Long-term records from systems aged 0--40 years confirm that most residual salt is removed in less than ten years, with effluent-salinity--drain-age relationships stable across decades \parencite{JOHNSTON1993}. Modeling of Haryana sites indicates 15--50 years to reach equilibrium effluent salinity, so drains may yield poor-quality water during much of their service life; drainage-water salinity is more sensitive to leaching fraction than to drain spacing, and raising the leaching fraction from 0.2 to 0.4 halved the equilibration time \parencite{Kelleners2000}. Field observations show drainage-water salinity first rising (initial mobilization of stored salt) then declining with a lag of about two years \parencite{Manjunatha2004}, decreasing to a mean of 2.6~dS~m$^{-1}$ suitable for winter irrigation \parencite{Sharma2006}, and progressively falling across cropping seasons as profile salinity declined \parencite{Jayawardane2001}. A cumulative-drainage relationship governs leaching: significant linear relations link cumulative drainage water to salt leached (R$^2=0.970$) and an exponential relation links drainage/soil depth to remaining/initial salt (R$^2=0.764$), with about 70\% of salt at 1.0~m depth removed by 339~mm of drainage over three years \parencite{Bahçeci2009}. Salt-mass export is strongly conditioned by geometry and geochemistry: drainage-water salinity and cumulative salt load increase with both drain depth and spacing, depth having the larger effect where groundwater quality declines with depth (an 80~m/4~m system exported roughly 36\% more than a 20~m system) \parencite{GRISMER1993}; deeper impervious layers and aquifer layering raise effluent salinity \parencite{KAMRA1991,KAMRA1994}. At watershed scale, region-wide tile drainage increased tile salt loading from 6 to 359~kg~ha$^{-1}$ while reducing groundwater salt loading, giving a modest net increase in total salt yield (1619 to 1651~kg~ha$^{-1}$) and a 12\% decline in aquifer SO$_4$ concentration, but producing episodic loading pulses that doubled or tripled baseline in-stream loads \parencite{Addab2022}. Nile Delta monitoring shows drainage-water salinity varying strongly in space and time, with secondary drains up to four times more saline than collectors \parencite{ElAgha2020}. The input/output salt-mass ratio stabilizes near 0.4, and higher application efficiency (lower leaching fraction) enhances salt buildup \parencite{KARAJEH1994c}.

\subsection{Vertical Redistribution and Ion-Specific Leaching}

Leaching under subsurface drainage is vertically stratified and ionically selective, a structure with direct implications for where and how fast reclamation occurs. Desalinization is strongest at the surface and weakens with depth: upper $>$ middle $>$ lower layer ordering is reported for continuous-head leaching \parencite{Chen2022}, deeper layers may resist or reverse (negative desalination in the 60--80~cm layer even as 0--20~cm exceeds 80\%) \parencite{Han2023}, and reduction magnitude is consistently greater in surface than subsoil layers \parencite{Narjary2024,Tian2024}. Soil texture modulates this: permeable sandy-loam and loam profiles sustain EC reduction while clayey sublayers limit leaching and can show negative EC change at depth \parencite{Qin2025}. Ion mobility follows characteristic sequences. Under heavy rainfall the order was Na$^+$ $>$ Cl$^-$ $>$ SO$_4^{2-}$ $>$ Mg$^{2+}$ \parencite{Yu2016}, and leaching efficiency was highest for Cl and Na (2018 leaching ratios of 91\% and 88\% at 0--20~cm), intermediate and irrigation-sensitive for Mg and SO$_4$, and lowest for Ca, which persists largely as precipitated gypsum ($>$90\% of surface Ca), whereas $>$90\% of Na is free Na$^+$ \parencite{Qian2025}. The chemical composition of drainage water resembles groundwater more than the leached-out salt \parencite{Qian2025}, and drainage significantly reduced chloride and sulfate in the profile \parencite{Sharma2000}.

A frequently observed complication is re-accumulation. Seasonal dynamics follow a three-phase ``wash-out--recovery--accumulation'' pattern, with re-accumulation positively correlated with reproductive growth (r~$=0.83$), Cl$^-$ and Na$^+$ leached downward, SO$_4^{2-}$ retained in the middle layer, and HCO$_3^-$ dominant near the surface; some wide-spacing treatments showed a 5--8~g~kg$^{-1}$ mid-layer increase signaling secondary-salinization risk \parencite{Zhu2025b}. Late-stage re-salinization occurred despite effective spring leaching, more severely under controlled than free drainage \parencite{Dou2022}, and surface salt accumulation reappears during dry, low-drainage periods \parencite{Yu2016,Shi2021b}. The direction of salt movement itself can be reoriented toward the drain, changing from deeper-zone or lateral movement to flow into the open drainage \parencite{Okuda2020}. Salinity distribution around drains is spatially organized: desalinization zones form within about 0.0--0.3~m of the drain \parencite{Lu2019}, the desalination rate is greatest immediately adjacent to the pipe and smallest farthest from it \parencite{Heng2021}, salinity increases with distance from the outlet \parencite{Han2023}, and normalized-salinity analysis places low-salinity zones within about 5~m and high-salinity zones beyond 7--8~m of the drain, with the greatest reduction within 7~m of a deep drain \parencite{Mirlas2003,Benyamini2005}.

\subsection{Leaching Hydraulics and Efficiency Enhancement}

A mechanistic body of work explains why leaching through subsurface drains is often inefficient and how efficiency can be raised, quantifying the non-uniform flow field between drains. Under complete ponding, flow concentrates near the drain: roughly 80\% of flow passes through 20\% of the soil volume, about 40\% of water after 450~h does not contribute to leaching, and desalinizing the full profile required 5400~h versus only 900~h for the top 1~m \parencite{RAO1991}. The flow fractions from successive strips fall steeply from the drain (0.819, 0.134, 0.035, 0.012) \parencite{Youngs2000}. Two families of remedies follow. Progressive/partial flooding leaches the central (poorly flushed) area first: sequential ponding achieved equal desalinization with about 50\% less water \parencite{RAO1991}, and a progressive four-strip schedule used only 1.7\% of the water and 8.3\% of the time of complete flooding (16.4\% water and 25.0\% time for infinitely deep soil), while also cutting evaporation \parencite{Youngs2000}. Low-permeability barriers or surface mulches over the drain redistribute flow toward uniform downward leaching: a barrier reduces wasted near-drain flow, with a 10~m barrier saving 52--82\% of leaching water and a half-extent barrier optimal \parencite{Yu2018}; a low-permeability surface mulch cut leaching timeframe and effluent volume by 64\% and 81\% and raised leaching efficiency, outperforming a buried barrier \parencite{Zhang2022}.

Soil structure and heterogeneity strongly govern efficiency. Macro-pores bypass surface infiltration and weaken the gradient toward the drain, increasing the time to leach to 0.6~m about twelvefold and water use roughly eighteenfold, with the effect worsening at greater macro-pore depth and density \parencite{Xin2016}. Layering can improve leaching relative to uniform soil, yet macro-pores still raise water use by more than an order of magnitude \parencite{Xin2016}. Mass-transfer limitations in dual-domain soils extend the leaching timeframe as trapped salt slowly diffuses from immobile domains; lower mass-transfer coefficient, larger capacity ratio, lower conductivity, wider spacing, and shallower depth all prolong leaching (immobile-domain leaching took 186 days versus 58 for the mobile domain in one case) \parencite{Zhang2024}. Salt movement in these systems is largely advection-dominated \parencite{Yu2018}. Applied-water strategy matters: at equal cumulative leaching water, multiple low-ponding cycles matched fewer high-ponding cycles, desalination increased with repeated cycles (about 64\% per saline-saturation drainage, reaching about 96\% after the fourth cycle), and deeper ponding (300~mm) improved desalting relative to shallow ponding \parencite{Jia2023}; deepening the pipe improved desalination more than widening the rice-husk zone \parencite{Jia2023}. More salt discharged under flood than drip leaching at the same discharge rate, and salt load rose linearly with discharge (R$^2>0.99$) \parencite{Li2022}. In fine-textured soils leaching efficiency is intrinsically low and declines with time, making salt removal a slow process \parencite{MANN1982}. Ditching and deep vertical rotary tillage enhance infiltration and leaching: ditching raised drainage per unit area by 14.65\% and surface (0--20~cm) desalination by 13.40\% where burial depth was under 1.3~m \parencite{Tian2023}, and deep vertical rotary tillage lowered bulk density, enhanced infiltration, and improved desalinization and yield \parencite{Li2026}.

\subsection{Saline and Brackish Water Reuse under Drainage}

A consolidated line of evidence establishes that subsurface drainage enables the safe agricultural reuse of saline and brackish water, provided salt balance is maintained by adequate leaching. Blended saline drainage water up to 9~dS~m$^{-1}$ produced about 90\% wheat yield in semi-arid Indian soils, with cyclic canal/drainage use generally outperforming continuous drainage-water use and long-term use causing only slight soil changes after monsoon leaching \parencite{Sharma2004}. Five-year data show relative wheat yields of 95.8\%, 90.3\%, 83.7\%, and 77.8\% at 6, 9, 12, and 18.8~dS~m$^{-1}$, with monsoon rain and pre-sowing irrigation keeping post-leaching root-zone salinity below 4~dS~m$^{-1}$ and no significant adverse effect on hydraulic conductivity \parencite{Sharma1998}. Conjunctive-use economics favor blending: mixing drainage water (12~dS~m$^{-1}$) with canal water optimized wheat yield at 4~dS~m$^{-1}$ and raised farm income \parencite{AGNIHOTRI1992}. In China, brackish-water leaching under drainage discharged salt exceeding the salt imported in favorable treatments (26.2--153.8~g replacement), with burial depth and irrigation amount controlling replacement and brackish-water salinity having little effect \parencite{Li2023}; saline (6~g~L$^{-1}$) irrigation matched freshwater wheat yields with no annual salt accumulation in the 0--45~cm layer and continued year-on-year salinity decline \parencite{Gao2023}; and drainage kept saline irrigation viable up to 6.25~dS~m$^{-1}$ for maize at 0.8~m depth \parencite{Feng2017}, 4.0--4.4~dS~m$^{-1}$ at 80~cm depth for grain \parencite{Feng2021}, and 8.0~dS~m$^{-1}$ under cut-soiler drainage without extra salt loading \parencite{Neha2022}. Optimal cotton outcomes occurred at brackish salinity of 2--4~g~L$^{-1}$ with 30~m spacing, the preferred depth increasing with salinity \parencite{Li2024}, and maximum acceptable seawater-mixture rates rose with leaching amount (37.27--52.57\% for 80--160~mm) \parencite{Lu2019}. Reuse nonetheless carries sodicity and accumulation risks: cyclic reuse supplying 70\% of summer irrigation caused profile salinity fluctuations and sodic surface soils (SARe~$>$~11), with EC rising from 2 to 6~dS~m$^{-1}$ despite winter-rain leaching \parencite{Mitchell2000}, and drainage-water SAR exceeded effluent SAR during reuse \parencite{Jayawardane2001}. Modeled reuse of poor-quality drain water over subsurface drains caused only marginal deterioration of shallow groundwater \parencite{KARAJEH1994c}. Blending criteria and rice-based reclamation for saline water are codified for the Nile Delta \parencite{Kotb2000}, and drainage/reduced-evaporation/low-EC-water combinations are highlighted as core salinity-control components in a global SICS review \parencite{Cuevas2019}.

\subsection{Integration with Complementary Practices}

Subsurface drainage performs best when integrated with irrigation, tillage, cropping, and hydraulic complements, and several such combinations are well documented. Controlled drainage conserves water and raises yield but trades off leaching: it reduced shallow ditch drainage by about 60\% and buffered salinity via stable deep drainage \parencite{Jia2007}, reduced subsurface drainage from 44.5 to 16.8~cm and the leaching factor to about 0.3 \parencite{Jia2007}, and improved crop water productivity while increasing late-stage re-salinization, with a 50~cm growth-stage depth optimal for sunflower \parencite{Dou2022}. Combined subsurface-pipe-plus-open-ditch systems desalinated effectively (drainage volume ratio about 9:1 favoring ditches), reaching 170\% of a single ditch's efficiency, with optimal pipe depth 90--110~cm, ditch depth 180--210~cm, and collection length 750--850~m \parencite{Yao2024}. Timing drainage to snowmelt and the rainy season kept the water table within 1~m, cut waterlogging by over two-thirds, more than doubled surface salt leaching, and produced residual benefits beyond active drainage periods \parencite{Yu2016}. Autumn/pre-planting irrigation coupled with drainage optimized emergence and raised sunflower yield 9.68--14.70\%, with a recommended 180~mm autumn quota \parencite{Liu2023,Liu2025c}. In cold semi-arid Canada, tall-wheatgrass windbreaks increased snow water equivalent (72.8 versus 25.5~mm) and, combined with drainage, gradually reduced root-zone salinity and multiplied barley yield, while pre-seeding irrigation lifted alfalfa emergence \parencite{Steppuhn2006,Steppuhn2017}. Cut-soiler-based preferential shallow drainage cut salinity 50--60\% over two years and raised pearl-millet and mustard yields, with drainage and salt removal greatest in the rainy season \parencite{Neha2022,Neha2022b}. Gypsum and cover-crop amendment improved soil physical properties and reduced crust strength under saline reuse, though gypsum did not lower EC \parencite{MANN1982,Mitchell2000}, and balanced N--K fertilization (150~kg~N and 50~kg~K~ha$^{-1}$) optimized wheat yield as drainage progressively lowered salinity \parencite{Singh2001b}. Where upward hydraulic pressure or perched conditions dominate, drainage design must adapt: combined drainage with relief wells was the most effective configuration for reducing salinity and lowering the water table within 2--3 years, whereas shallow drainage alone failed under upward pressure \parencite{Benyamini2005,Mirlas2003}. Low-cost alternative technologies---straw-drainage modules performing comparably to conventional pipe systems \parencite{Lu2019} and cut-drain shallow drainage that reduced remaining salt to a third or quarter of the control and raised cotton yield about 18\% \parencite{Okuda2020}---extend the approach to resource-limited settings. Regional coupling of optimized drip irrigation with tailored subsurface drainage reduced district-wide irrigation by 10\% and soil salinity by 15\% while raising yields \parencite{Liu2026}.

\subsection{Modeling, Prediction, and Monitoring Tools}

The quantitative understanding above rests on a mature toolkit of validated models and monitoring methods, which is itself a consolidated finding. Physically based drainage--salinity models reproduce field water tables and salt dynamics well: DRAINMOD and DRAINMOD-S captured water-table fluctuations (NSE up to 0.84) and soil salinity (NSE 0.40--0.92) and were used to derive spacing/depth recommendations \parencite{Chen2025,Tian2023,Tian2023c,Tian2024,Davoodi2019,Li2024}. HYDRUS variants (1D/2D/3D) and their couplings simulated water--salt transport, drainage mass, and crop response accurately across brackish-leaching, forest, tillage, and controlled-drainage settings \parencite{Li2023,Dou2022,Wang2024c,Li2026,Jia2023,Liu2022}, and coupled HYDRUS--EPIC and AquaCrop/DRAINMOD-S frameworks reliably reproduced soil water--salt, yield, and evapotranspiration for evaluating long-term saline-water irrigation \parencite{Feng2019,Feng2021,Li2024}. Salt-and-water-balance models---SaltMod, SahysMod, SWMS-2D, and SALTIRSOIL\_M---predicted root-zone salinity, drain discharge, and effluent salinity closely enough for management use \parencite{Bahçeci2006,Bahçeci2007,Yao2017,Qian2021b,Abdelkrim2021}, and finite-element and analytical solutions for reactive-solute transport, dual-domain mass transfer, and barrier/mulch schemes matched observations and yielded design guidelines \parencite{KAMRA1991,KAMRA1994,Zhang2022,Zhang2024,Youngs2000}. Watershed-scale SWAT-Salt tracked ion loading through tile, groundwater, and runoff pathways and linked local tile behavior to downstream salinity \parencite{Addab2022,Liu2026}, while long-validated empirical effluent-quality relationships remain reliable across decades \parencite{JOHNSTON1993,Kelleners2000}. On the monitoring side, electromagnetic-induction and apparent-EC surveys with geostatistical and time-specific-regression methods enable rapid, nondestructive, depth-resolved assessment of drainage-system efficacy and salinity mapping \parencite{Narjary2024,Song2023,Yu2025,Singh2020d,Mirlas2003}, and emerging IoT-plus-LSTM digital-twin systems achieved high-accuracy soil-EC and water-content forecasting (test R$^2$ up to 0.97) for coordinated irrigation--drainage control \parencite{Qin2025}. Comprehensive field characterizations of irrigation-induced salinity and shallow-groundwater conditions provide the baseline against which these tools are calibrated \parencite{Gates2002}. Collectively, this modeling and monitoring capability is what allows the design and salt-balance principles summarized above to be transferred with confidence to new sites and management scenarios.